\subsection{Conditional-tail information beyond conditional mean and variance}
\label{app:conditional-tail-information}
The following exact bounded-law example isolates information lost even
after conditional mean and variance have been obtained. It concerns the
lower-tail feasibility criterion, rather than unconditional dominance of
expected return. It is an analytical score witness, not a physical
measurement or an empirical calibration certificate.

\begin{proposition}[Equal conditional moments, different paid-tail actions]
\label{prop:conditional-tail-information}
At one fixed current source configuration and calibration context, let
two conditional laws be
\begin{equation}
 \begin{aligned}
 P(U=0)&=P(U=1/5)=1/2,\\
 Q(U=-1/5)&=1/10,\qquad Q(U=2/15)=9/10.
 \end{aligned}
 \label{eq:conditional-tail-witness-laws}
\end{equation}
Both have conditional mean $1/10$ and variance $1/100$, but
\begin{equation}
 \begin{aligned}
 \operatorname{LT}_{1/2}(P)&=0,&
 \operatorname{LT}_{1/2}(Q)&=1/15,\\
 128\operatorname{LT}_{1/2}(P)-5&=-5,&
 128\operatorname{LT}_{1/2}(Q)-5&=53/15.
 \end{aligned}
 \label{eq:conditional-tail-witness-actions}
\end{equation}
Consequently no policy using only this context and these conditional
moments can reproduce both prescribed positive-tail admission decisions.
Among such policies required to admit only when the true paid lower tail
is nonnegative for \emph{every} law compatible with those moments,
both configurations must be rejected. The full-law criterion admits under
$Q$, where its conditional expected conservative paid return is $39/5$.
This is the value of resolving tail feasibility; unconstrained admission
under both laws has the same $39/5$ expected paid return and is not
dominated in the mean-return criterion by this argument.
\end{proposition}
\begin{proof}
For $P$, the mean is $(1/2)(1/5)=1/10$ and the second moment is
$(1/2)(1/25)=1/50$. For $Q$, these are respectively
$-(1/10)(1/5)+(9/10)(2/15)=1/10$ and
$(1/10)(1/25)+(9/10)(4/225)=1/50$.
Subtracting the squared mean gives variance $1/100$ in each case.
The lowest half of $P$ is entirely at zero. For $Q$, its lowest half
contains mass $1/10$ at $-1/5$ and mass $2/5$ at $2/15$, giving
\begin{equation}
 \operatorname{LT}_{1/2}(Q)
 =2\left\{-\frac1{10}\frac15+\frac25\frac2{15}\right\}
 =\frac1{15}.
\end{equation}
Translation and positive scaling give the paid-tail scores in
Eq.~\eqref{eq:conditional-tail-witness-actions}. The conditional
expected paid return in either law is $128/10-5=39/5$.
A moment-only policy sees identical inputs and must use the same action
or admission probability. Since $P$ is compatible and has paid lower
tail $-5$, any positive admission probability violates the stated
law-by-law tail-feasibility requirement. Rejection under both laws
therefore misses the feasible $39/5$ expected action under $Q$.
\end{proof}
The example assumes a common zero score correction and a common feasible
execution ledger. It compares exact conditional distributions, so it does
not require a particular Bayesian or non-Bayesian estimator. These atoms
give fractional gross values at $N=128$; they illustrate the bounded-law
scoring functional rather than a literal enumeration of integer
correctness counts in the recorded service contract. The separate
physical tests must establish whether retained pairing changes such
tail decisions accurately and improves complete paid return.
